\documentclass[11pt]{article}
\usepackage[a4paper,margin=27mm]{geometry}
\usepackage{amsmath,amssymb,amsthm}
\usepackage[hidelinks]{hyperref}
\setlength{\emergencystretch}{3em}
\newtheorem{proposition}{Proposition}
\title{A counterexample to TLMC conjecture 00000002320}
\author{AI-assisted mathematical exploration}
\date{September 21, 2026}

\begin{document}
\maketitle

\begin{abstract}
The conjecture places the probabilities of generating all finite groups by two
uniform random elements in the interval $[1/4,1]$.  The elementary abelian group
$(\mathbb Z/2\mathbb Z)^3$ has no generating pair, so its probability is zero.
We describe the standard counting definition, the counterexample, and its
correspondence with the accompanying Lean source.
\end{abstract}

\section*{Target and conventions}
Conjecture 00000002320 of The Last Math Competition, at the fixed revision
identified in reference~\cite{tlmc}, asserts that the spectrum of generating-pair
probabilities over \emph{all finite groups} is a dense subsequence of $[1/4,1]$.
It also makes a separate assertion about $\operatorname{PSL}(2,p)$.
We disprove the necessary consequence
\[
  P(G)\geq \frac14\qquad\text{for every finite group }G.
\]
The source does not restrict the quantifier to simple groups or to groups that
can be generated by two elements.

For a finite group $G$, let
\[
 N_2(G)=\bigl|\{(a,b)\in G\times G:\langle a,b\rangle=G\}\bigr|,
 \qquad P(G)=\frac{N_2(G)}{|G|^2}.
\]
Thus the two elements are sampled independently and uniformly, with replacement;
pairs are ordered.  This is the standard definition of the two-generator
probability; see Pak~\cite[p.~3]{pak} for the general $k$-tuple definition.

\begin{proposition}
Let $G=(\mathbb Z/2\mathbb Z)^3$.  Then $|G|=8$ and $P(G)=0$.
Consequently the spectrum asserted in conjecture 00000002320 is not contained
in $[1/4,1]$.
\end{proposition}

\begin{proof}
Write the group operation additively.  For any $a,b\in G$, consider
\[
 H_{a,b}=\{0,a,b,a+b\}.
\]
The group is abelian and every element satisfies $x+x=0$.  It follows that
$H_{a,b}$ contains zero and is closed under addition and inverses, so it is a
subgroup containing $a$ and $b$.  Therefore
\[
  \langle a,b\rangle\subseteq H_{a,b},
  \qquad |\langle a,b\rangle|\leq 4<8=|G|.
\]
No pair generates $G$.  There are $8^2=64$ ordered pairs in total and zero
generating pairs, giving
\[
  P(G)=\frac{0}{64}=0<\frac14.
\]
This contradicts the required interval containment.
\end{proof}

\section*{Lean correspondence and verification scope}
\begingroup\raggedright
The accompanying file \texttt{Counterexample2320.lean} uses the actual mathlib
group \texttt{Multiplicative (ZMod 2 $\times$ ZMod 2 $\times$ ZMod 2)}.
The multiplicative type tag only changes the notation for the operation.
The generating relation is mathlib's \texttt{Subgroup.closure} of the two-element
set, equal to the top subgroup.

The declarations \texttt{fourSubgroup} and
\texttt{\detokenize{pair_closure_le_four}} construct a subgroup with carrier
$\{1,a,b,ab\}$ and contain the generated subgroup in it.
The finite theorem \texttt{\detokenize{four_misses_element}} supplies an element
outside this subgroup for every pair.  Consequently
\texttt{\detokenize{no_pair_generates}} proves that no generated subgroup is top.
The definition \texttt{generatingPairs} filters the full finite set of ordered
pairs by this exact generation predicate, and
\texttt{generatingPairProbability} is its cardinality divided by $|G|^2$.
The terminal declarations are
\texttt{\detokenize{generatingPairProbability_zero}} and
\texttt{\detokenize{conjectureInterval_false}}.
\par\endgroup

The probability is represented exactly as a rational number.  Its usual
embedding into the reals preserves the value zero and the comparison with
$1/4$.  Finite checks use kernel-checked \texttt{decide}; no
\texttt{\detokenize{native_decide}} is used.  Compilation, axiom auditing, and
replay evidence are recorded separately in the delivered validation logs.

The argument refutes one necessary consequence of the conjecture.  It does
not claim a formal proof concerning the separate $\operatorname{PSL}(2,p)$
limit.  The fact about elementary abelian groups is standard, and no novelty
in group theory is claimed.  A bounded check of the repository's 100 pull
request titles and bodies at 2026-09-20 17:13 UTC found no matching submission
identifier; this is not a claim of worldwide priority.

\begin{thebibliography}{9}
\raggedright
\bibitem{tlmc}
The Last Math Competition,
\emph{Conjecture 00000002320}, revision
\texttt{95acb520ec5607c826b8a997b1ef2fc82d6f7c57}.
\href{https://github.com/The-Last-Math-Competition/The-Last-Math-Competition/blob/95acb520ec5607c826b8a997b1ef2fc82d6f7c57/conjectures/00000002320.md}{Fixed upstream source}.

\bibitem{pak}
Igor Pak,
\emph{On probability of generating a finite group},
December 30, 1999, especially p.~3.
\url{https://www.math.ucla.edu/~pak/papers/sim.pdf}.
\end{thebibliography}
\end{document}
